\section{Two-variable logic and tree size}\label{sec:FOtwo}
The formula $G_k$ is \FOtwo-definable, but its displayed linear-size
form uses three variable names. This section quantifies the difference,
proves the term-size bound of Theorem~\ref{thm:tree}, and separates
\CoR\ from \FOtwo\ (\S\ref{sec:reverse}).

\subsection{Quantifiers and fold sensitivity}\label{sec:quantifier-cost}
Recall the linear translation $\tau$ from \FOtwo\ formulas into \CoR\
given in \S\ref{sec:FOtwo-translation}. On the reference structure,
the off-diagonal support of $R\cap\id$ is empty. Therefore neither
of the compositions used to represent a reflexive atom,
$(R\cap\id);\one$ or $\one;(R\cap\id)$, notices a fold
(see Lemma~\ref{lem:notice}). Every quantifier, by contrast, contributes
one \emph{one-sided} composition $X;\one$ or $\one;X$; by the same
lemma, it can notice at most one fold. Boolean operations notice none.
This gives a bound in terms of quantifiers, rather than all
composition occurrences in $\tau(\varphi)$.
\begin{lemma}[Quantifier cost]\label{lem:FO-cost}
On the reference structure, every \FOtwo\ formula $\varphi$ satisfies
\[
 \tcost(\tau(\varphi))\le \#_{\exists\forall}(\varphi).
\]
\end{lemma}
The bound holds without restricting the Boolean connectives or the
quantifier alternation of the target formula.

Combining Lemma~\ref{lem:FO-cost} with Theorem~\ref{thm:test-cost} and
the tests of Lemma~\ref{lem:G-tests} gives
\[
 \varphi\equiv G_k\Longrightarrow \#_{\exists\forall}(\varphi)\ge2^k.
\]
In fact, agreement with $G_k$ on all finite structures suffices for
this bound. The \FOtwo\ expansion obtained from the DNF of \S\ref{sec:main-results}
uses exactly $2^k$ quantifiers, so the optimum quantifier count for
$G_k$ is exactly $2^k$. No restriction to existential target formulas is
made.

\subsection{Counting input occurrences in trees}\label{sec:trees}
The test--cost theorem also yields a stronger bound for syntax trees.
A \defn{singleton selector} for a color $\lambda$ is a
Boolean formula in the tables of the input relations, with constants
allowed, whose support on the entire palette is exactly $\{\lambda\}$.
\begin{restatable}[Disjoint input-occurrence charging]{lemma}{treecharging}\label{lem:tree-charge}
Fix the reference structure and its folds. Suppose every singleton
selector requires at least $d>0$ input-relation occurrences.
Then every \CoR\ term $T$, whether correct or not, satisfies
\[
 \tcost(T)\le \relatoms{T}/d.
\]
\end{restatable}
Appendix~\ref{app:trees} proves the lemma by assigning disjoint sets of
input occurrences to composition nodes. The disjointness uses the tree
structure of a term.

For the input relations of \S\ref{sec:G}, a Boolean formula in their
tables that contains neither $P_i$ nor $Q_i$ takes the same value on a
color and on its $i$-th-bit flip. Hence a singleton selector contains
$P_i$ or $Q_i$ for every $i$, so $d=k$, and this is exact, since
$H_\lambda$ selects $\lambda$. Theorem~\ref{thm:test-cost} and
Lemma~\ref{lem:tree-charge} give $\relatoms{T}\ge k2^k$ for every \CoR\
term $T$ agreeing with $G_k$ on all finite structures. By the
tree-size inequality from \S\ref{sec:prelim-resources},
\[
 |T|\ge2k2^k-1\qquad(T\equiv G_k),
\]
which is the lower bound of the following theorem.
\begin{theorem}[Term size]\label{thm:tree}
For every $k\ge1$, every \CoR\ term $T\equiv G_k$ has at least $2k2^k-1$
nodes. The lower bound holds even if $T$ and $G_k$ agree only on
finite structures, and $G_k$ is equivalent on all nonempty structures
to a \CoR\ term with $(2k+6)2^k-5$ nodes.
\end{theorem}
\begin{proof}
The lower bound is the inequality above. Factoring the conditions on the
read edge into a binary decision tree gives a term with $(2k+6)2^k-5$
nodes that is equivalent to $G_k$ on all nonempty structures
(Appendix~\ref{app:factored}).
\end{proof}

\subsection{Sentences and the variable hierarchy}\label{sec:sentences}
The preceding subsections establish the binary-query bounds of
Theorem~\ref{cor:sentences}. We now prove its sentence extension for
$G_k^\forall=\forall x\,\forall y\,(x=y\lor G_k(x,y))$.

\sentencebounds*
\paragraph{Quantitative and sentence extensions.}
\nakamura{sentenceの場合もステートメントにしてはどうか}
We also prove two refinements. Every \FOtwo\ formula defining $G_k$
has at least $k2^k$ relation-atom occurrences, and its minimum
syntax-tree size is $\Theta(k2^k)$. Moreover, $G_k^\forall$ is an
\FOtwo-definable \FOthree\ sentence of size $O(k)$ and quantifier
depth three, but every equivalent \FOtwo\ sentence over finite
structures has at least $2^k$ quantifiers and $k2^k$ relation-atom
occurrences; its minimum \FOtwo\ tree size is $\Theta(k2^k)$.
\begin{proof}[Proof of the quantitative and sentence extensions]
The exact $2^k$-quantifier lower bound for the binary query $G_k$
was established in \S\ref{sec:quantifier-cost}. For any equivalent \FOtwo\
formula $\varphi$, apply Theorem~\ref{thm:test-cost} and
Lemma~\ref{lem:tree-charge} to $\tau(\varphi)$. By the
relation-atom preservation established in \S\ref{sec:FOtwo-translation},
this gives
$\relatoms{\varphi}\ge k2^k$. The \FOtwo\ DNF expansion from
\S\ref{sec:main-results} has $O(k2^k)$ nodes and exactly $2^k$
quantifiers, proving the tree-size claim for the binary query.

For the sentence $G_k^\forall$, the source has size $8k+4$ and
quantifier depth three. By Lemma~\ref{lem:G-tests}, it is true in
$\mathcal M_0$ and false in every $\mathcal M_\lambda$. On each of
these structures its assignment relation is $\zero$ or $\one$, and
both relations are fixed by prediction. Thus these structures give $2^k$ tests for the
sentence itself. For any equivalent \FOtwo\ sentence $\varphi$,
apply Theorem~\ref{thm:test-cost} to $\tau(\varphi)$.
Lemma~\ref{lem:FO-cost} gives $\#_{\exists\forall}(\varphi)\ge2^k$,
and Lemma~\ref{lem:tree-charge}, together with exact preservation of
$\#_{\mathrm{rel}}$, gives $\relatoms{\varphi}\ge k2^k$.
Conversely, substitute the \FOtwo\ expansion of $G_k$ into the
sentence. It has $2^k+2$ quantifiers and $O(k2^k)$ nodes, proving the
sentence upper bounds and the matching order of tree size.
\end{proof}
The quantifier optimum for this sentence is therefore in
$[2^k,2^k+2]$; the exact $2^k$ optimum proved earlier concerns the binary
query.

This complements the variable-saving separations of
Vinall-Smeeth~\citep[Theorem~2]{vinall-smeeth2024}, whose stated range starts at three
variables on the target side. Our result concerns the remaining comparison
of three and two variables, over growing binary signatures. It does not
claim a lower bound for separating a single prescribed pair of structures
as a function of their size. On pure finite linear orders, in contrast,
Grohe and Schweikardt~\citep[Theorem~4.3]{grohe-schweikardt2005} give a polynomial translation
from three-variable sentences to two-variable sentences. The qualifications
``pure orders'' and ``sentences'' are essential to that comparison.

\subsection{Separating \texorpdfstring{{\boldmath\CoR}}{CoR} from
 \texorpdfstring{{\boldmath\FOtwo}}{FO2}}\label{sec:reverse}
The linear translation of \S\ref{sec:FOtwo-translation} goes from \FOtwo\ to \CoR. On the reference
structure, a composition notices at most two folds and a quantifier at
most one, so tests on the reference structure bound \CoR\ circuits and
\FOtwo\ formulas alike, up to a factor of two. To separate the two
formalisms, we use a structure on which a quantifier still notices at
most one fold, but a single composition notices all of them.

Use input symbols $R,P_1,\ldots,P_k$, and put
\[
 \mathscr P_i^1\coloneqq\id\cap P_i,\qquad \mathscr P_i^0\coloneqq\id\cap\overline{P_i},\qquad
 U_k\coloneqq R;\bigcap_{i=1}^k
 \left[(\mathscr P_i^0;\one;\mathscr P_i^0)\cup(\mathscr P_i^1;\one;\mathscr P_i^1)\right].
\]
The complement in $\mathscr P_i^0$ is taken before intersection with the identity;
it is not the complement of $\mathscr P_i^1$. Expanding these abbreviations gives
an $O(k)$-size tree. For a vertex $v$, let
$\pi(v)\coloneqq\{i\in[k]:P_i(v,v)\}$ be its \defn{loop pattern}.
Here $\pi(v)$ is a local one-vertex notation, distinct from the pair
profile $\pi(u,v)$ of \S\ref{sec:domain-size}.
Then $U_k(x,y)$ holds exactly when some $R$-successor $z$ of $x$
satisfies $\pi(z)=\pi(y)$; equivalently, $U_k$ is defined by the
\FOthree\ formula
\[
 \exists z\left(R(x,z)\land
     \bigwedge_{i=1}^k(P_i(z,z)\leftrightarrow P_i(y,y))\right).
\]

\paragraph{The structures.}
Let $\mathbb Z_{11}$ denote the integers modulo $11$, and put
$V\coloneqq2^{[k]}\times\mathbb Z_{11}$ and
$\Gamma\coloneqq\{1,2,3,4,7,8\}\subseteq\mathbb Z_{11}$. The structure
$\mathcal W_0$ on $V$ has the loops $P_i((A,s),(A,s))$ for $i\in A$, no
other $P_i$-entries, and $R((A,s),(B,s'))$ exactly when $s'-s\in\Gamma$;
thus $\pi((A,s))=A$, and $R$ has no loops. The \defn{atomic type} of a pair
$(u,v)$ of distinct vertices is $(R(u,v),R(v,u),\pi(u),\pi(v))$, and that
of a loop $(u,u)$ is $\pi(u)$; it determines the values of all atomic
formulas at the pair. As $s'-s$ ranges over
$\mathbb Z_{11}\setminus\{0\}$, the truth values of $s'-s\in\Gamma$ and
$s-s'\in\Gamma$ take each of the four combinations at least twice:
$(1,0)$ at $s'-s\in\{1,2\}$, $(0,1)$ at $s'-s\in\{9,10\}$, $(1,1)$ at
$s'-s\in\{3,4,7,8\}$, and $(0,0)$ at $s'-s\in\{5,6\}$. Hence every
vertex $u$ is the first entry of at least two off-diagonal pairs of each
atomic type $(r,r',\pi(u),B)$, and the second entry of at least two
of each type $(r,r',B,\pi(u))$, for $r,r'\in\{0,1\}$ and $B\subseteq[k]$.

Fix the hub $h\coloneqq(\varnothing,0)$. For $A\subseteq[k]$, the fold
of $A$ yields the structure $\mathcal W_A$ obtained from $\mathcal W_0$ by
removing from $R$ the edges $(h,v)$ with $v\ne h$ and $\pi(v)=A$. It
changes the atomic types of these pairs and of their converses, and of no
other pair. Afterwards, no pair with first entry $h$ has a type
$(1,r',\varnothing,A)$, and no pair with second entry $h$ has a type
$(r,1,A,\varnothing)$; all other types still occur at $h$.

\paragraph{Type tables and predictions.}
As in \S\ref{sec:types}, we use relations determined by a finite
set of pair types, but here the types are atomic pair types rather
than edge colors, and the folds delete edges rather than recolor them.
A relation $X\subseteq V^2$ has a \defn{type table} if $X(u,v)$ depends
only on the atomic type of $(u,v)$ in $\mathcal W_0$; let $S_X$ be the
set of atomic types at which $X$ holds, rather than a set of colors.
Every atomic type of $\mathcal W_A$ occurs in $\mathcal W_0$, so the prediction $\pred XA$ can
read the same table at the atomic types of $\mathcal W_A$. As in
\S\ref{sec:unnoticed}, for operand values $\mathbf U$ having type tables,
an operation $\operatorname{op}$ whose result $\operatorname{op}(\mathbf U)$
also has a type table notices the fold of $A$ if
$\operatorname{op}(\pred{\mathbf U}A)\ne\pred{\operatorname{op}(\mathbf U)}A$,
and Boolean operations never notice a fold.

For each \FOtwo\ formula $\psi$, write $X_\psi$ for its binary
assignment relation in the structure under consideration, as in
\S\ref{sec:prelim-logic}. Existential quantifiers act as the projections
\[
 X_{\exists y\,\psi}=X_\psi;\one,\qquad
 X_{\exists x\,\psi}=\one;X_\psi.
\]
Universal quantifiers are handled by complement.
By induction on the formula, every relation that an \FOtwo\ formula
defines on $\mathcal W_0$ has a type table: if $X$ has one, then
$(X;\one)(u,v)$ holds exactly when $S_X$ contains the atomic type of some
pair with first entry $u$, and these types are determined by $\pi(u)$;
symmetrically for $\one;X$.

\begin{lemma}[Projections notice at most one fold]\label{lem:W-projection}
If $X$ has a type table, then $X;\one$ and $\one;X$ each notice the fold
of at most one $A\subseteq[k]$.
\end{lemma}
\begin{proof}
At a pair $(u,v)$, the prediction of $X;\one$ holds exactly when $S_X$
contains the atomic type in $\mathcal W_0$ of some pair with first entry
$u$, and $\pred XA;\one$ holds exactly when $S_X$ contains the atomic type
in $\mathcal W_A$ of some such pair. If $u\ne h$, only the
off-diagonal pair $(u,h)$ can change. Every off-diagonal atomic type
realized among pairs with first entry $u$ is realized at least twice,
so its old type is still realized after the fold. Its new type was
already realized in $\mathcal W_0$, while the diagonal pair $(u,u)$
is unchanged. Hence the sets of realized types coincide. If $u=h$, the fold removes exactly the types
$(1,r',\varnothing,A)$ and adds none. Hence $X;\one$ notices the fold of
$A$ only if the types in $S_X$ of pairs with first entry $h$ are all of
the form $(1,r',\varnothing,A)$, and there is at least one; this
determines $A$. The argument for $\one;X$ uses second entries.
\end{proof}

\reverse*
\paragraph{Equational sentence extension.}
The \CoR\ equation $\id\cup U_k=\one$ is \FOtwo-definable as a
sentence, but every equivalent \FOtwo\ sentence over finite
structures has at least $2^k$ quantifiers.
\begin{proof}[Proof of the query and sentence claims]
For the upper bounds, let $\vartheta_A(u)$ be the quantifier-free
formula asserting that $\pi(u)=A$. The \FOtwo\ formula
\[
 \bigvee_{A\subseteq[k]}\left[\vartheta_A(y)\land
       \exists y\,(R(x,y)\land\vartheta_A(y))\right]
\]
defines $U_k$ with exactly $2^k$ quantifiers; the inner $y$ is rebound
only in the second conjunct. The equation $\id\cup U_k=\one$ says that
for all $x$ and all $y\ne x$, some $R$-successor of $x$ has the loop
pattern of $y$, and so does the \FOtwo\ sentence
\[
 \forall x\bigwedge_{A\subseteq[k]}\left[
 \exists y\,(x \neq y\land\vartheta_A(y))\rightarrow
 \exists y\,(R(x,y)\land\vartheta_A(y))\right].
\]

For the lower bounds, every vertex of $\mathcal W_0$ has an $R$-successor
with each loop pattern, so $U_k$ is the full relation on $\mathcal W_0$
and the equation holds there. On $\mathcal W_A$, the hub has no
$R$-successor with loop pattern $A$, so $U_k(h,v)$ fails for every
$v\ne h$ with $\pi(v)=A$, and the equation fails. Let $\psi$ be an
\FOtwo\ formula agreeing with $U_k$ on all finite structures, or an
\FOtwo\ sentence holding on exactly the finite structures satisfying
the equation. Its value on $\mathcal W_0$ is the full relation, whose
prediction on $\mathcal W_A$ is again full, but its value on
$\mathcal W_A$ is not. If no quantifier of $\psi$ noticed the fold of
$A$, then, as in the proof of Lemma~\ref{lem:unnoticed}, every subformula
would take its predicted value on $\mathcal W_A$, by induction on the
formula: atomic types determine atomic formulas, and Boolean operations
notice nothing. Hence some quantifier notices the fold of each $A$, and
by Lemma~\ref{lem:W-projection} each quantifier notices at most one, so
$\psi$ has at least $2^k$ quantifiers.
\end{proof}
